% ---------------------------------------------------------------
%  Page geometry and typography
% ---------------------------------------------------------------
\usepackage[paperwidth=7.5in,paperheight=9.75in,
            inner=0.95in,outer=0.78in,top=0.82in,bottom=0.88in,
            headsep=16pt,footskip=28pt]{geometry}
\usepackage[T1]{fontenc}
\usepackage{lmodern}
\usepackage{microtype}
\usepackage{amsmath,amssymb,mathtools}
\usepackage{graphicx}
\usepackage{tikz}
\usetikzlibrary{arrows.meta,calc,angles,quotes,patterns,decorations.pathmorphing,
                decorations.markings,intersections,positioning,shapes.geometric,bending}
\usepackage{pgfplots}
\pgfplotsset{compat=1.18}
\usepackage{caption}
\usepackage{enumitem}
\usepackage{fancyhdr}
\usepackage{titlesec}
\usepackage{xcolor}
\usepackage{booktabs}
\usepackage{multirow}
\usepackage{multicol}
\usepackage{siunitx}
\sisetup{detect-all,per-mode=symbol}
\usepackage[hidelinks]{hyperref}
\hypersetup{
  pdftitle={Training and Selection for the International Physics Olympiad},
  pdfauthor={Zheng Yongling (Editor-in-Chief)},
  pdfsubject={Physics Olympiad problems, solutions and experiments},
  pdfkeywords={physics, olympiad, mechanics, thermodynamics, electromagnetism, optics, atomic physics}
}

\definecolor{bookblue}{RGB}{20,45,85}

% ---------------------------------------------------------------
%  Running heads
% ---------------------------------------------------------------
\pagestyle{fancy}
\fancyhf{}
\fancyhead[LE]{\small\scshape\nouppercase{\leftmark}}
\fancyhead[RO]{\small\scshape\nouppercase{\leftmark}}
\fancyfoot[LE,RO]{\small\thepage}
\renewcommand{\headrulewidth}{0.4pt}
\fancypagestyle{plain}{%
  \fancyhf{}%
  \fancyfoot[LE,RO]{\small\thepage}%
  \renewcommand{\headrulewidth}{0pt}}

% ---------------------------------------------------------------
%  Chapter and part styles
% ---------------------------------------------------------------
\titleformat{\chapter}[display]
  {\normalfont\raggedright}
  {\color{bookblue}\scshape\large\chaptertitlename\ \thechapter}
  {10pt}
  {\color{bookblue}\Huge\bfseries}
\titlespacing*{\chapter}{0pt}{18pt}{36pt}

\titleformat{\section}
  {\normalfont\large\bfseries\color{bookblue}}{\thesection}{0.7em}{}
\titleformat{\subsection}
  {\normalfont\normalsize\bfseries}{\thesubsection}{0.6em}{}

\renewcommand{\partname}{}
\titleformat{\part}[display]
  {\centering\normalfont}
  {}
  {0pt}
  {\Huge\scshape\color{bookblue}}
\titlespacing*{\part}{0pt}{60pt}{50pt}

% ---------------------------------------------------------------
%  Environments for problems, hints and solutions
% ---------------------------------------------------------------
\newenvironment{problem}[1]
  {\par\addvspace{9pt}\noindent\textbf{Problem~#1.}\hspace{0.55em}\ignorespaces}
  {\par\addvspace{3pt}}

\newenvironment{hint}[1]
  {\par\addvspace{8pt}\noindent\textbf{Hint~#1.}\hspace{0.55em}\ignorespaces}
  {\par\addvspace{3pt}}

\newenvironment{solution}[1]
  {\par\addvspace{10pt}\noindent\textbf{Solution~#1.}\hspace{0.55em}\ignorespaces}
  {\par\addvspace{5pt}}

\newenvironment{note}
  {\par\addvspace{6pt}\begin{small}\noindent\textit{Note.}\hspace{0.5em}\ignorespaces}
  {\end{small}\par\addvspace{4pt}}

\newenvironment{commentary}
  {\par\addvspace{5pt}\begin{small}\noindent\textit{Comment.}\hspace{0.5em}\ignorespaces}
  {\end{small}\par\addvspace{4pt}}

% Figure box: usage \begin{figbox}{1.28} ... tikzpicture ... \end{figbox}
\newenvironment{figbox}[1]
  {\par\addvspace{6pt}\def\figboxcaption{#1}\begin{center}%
   \begin{minipage}{0.96\linewidth}\centering}
  {\par\smallskip{\small Fig.~\figboxcaption}\end{minipage}\end{center}\addvspace{5pt}}

\captionsetup{font=small,labelfont=bf,skip=4pt}

% ---------------------------------------------------------------
%  Mathematics shorthands
% ---------------------------------------------------------------
\newcommand{\dd}{\mathrm{d}}
\DeclareMathOperator{\arccot}{arccot}
\DeclareMathOperator{\arcsec}{arcsec}
\DeclareMathOperator{\arccsc}{arccsc}
\newcommand{\ee}{\mathrm{e}}
\newcommand{\ii}{\mathrm{i}}
\newcommand{\degree}{^{\circ}}
\newcommand{\avg}[1]{\overline{#1}}
\newcommand{\unitv}[1]{\,\mathrm{#1}}
\newcommand{\dg}{\degree}
\newcommand{\vs}{\vspace{2pt}}
\setlength{\parindent}{1.2em}
\setcounter{tocdepth}{2}
\setlength{\parskip}{0pt}
\emergencystretch=1.2em
\clubpenalty=1000
\widowpenalty=1000

% Standard TikZ styles used throughout the book
\tikzset{
  >={Stealth[length=1.8mm,width=1.3mm]},
  axis/.style={->,line width=0.5pt},
  thinline/.style={line width=0.5pt},
  thickline/.style={line width=0.9pt},
  dashedline/.style={dashed,line width=0.5pt},
  force/.style={->,line width=0.7pt},
  dim/.style={<->,line width=0.4pt},
  hatch/.style={pattern=north east lines,pattern color=black!55},
  ball/.style={circle,draw,line width=0.5pt,inner sep=0pt,minimum size=5pt},
  mass/.style={rectangle,draw,line width=0.5pt,inner sep=1.5pt},
  spring/.style={decorate,decoration={zigzag,segment length=2.2mm,amplitude=1.1mm}}
}

\usepackage{makeidx}
\makeindex
